\documentclass[11pt]{article}

\usepackage[T1]{fontenc}
\usepackage{lmodern}
\usepackage{amsmath,amssymb,amsthm,mathtools}
\usepackage{microtype}
\usepackage[margin=1in]{geometry}
\usepackage{booktabs}
\usepackage{longtable}
\usepackage{enumitem}
\usepackage{xcolor}
\usepackage{listings}
\usepackage{hyperref}
\usepackage{url}

\hypersetup{
  colorlinks=true,
  linkcolor=blue!45!black,
  urlcolor=blue!55!black,
  citecolor=blue!45!black,
  pdftitle={BusyLean Contender 7: Autonomous Critical-Point Reflection},
  pdfauthor={OpenAI research assistant for Vladimir Reshetnikov}
}

\lstdefinelanguage{Lean}{
  keywords={def,theorem,inductive,namespace,end,open,import,where,match,with,fun,let,in,by,exact,change,forall,termination_by},
  sensitive=true,
  morecomment=[l]{--},
  morecomment=[s]{/-}{-/},
  morestring=[b]"
}
\lstset{
  language=Lean,
  basicstyle=\ttfamily\small,
  keywordstyle=\bfseries\color{blue!45!black},
  commentstyle=\itshape\color{black!58},
  frame=single,
  framerule=0.35pt,
  rulecolor=\color{black!30},
  columns=fullflexible,
  keepspaces=true,
  breaklines=true,
  showstringspaces=false,
  xleftmargin=0.75em,
  xrightmargin=0.75em
}

\newtheorem{theorem}{Theorem}
\newtheorem{lemma}[theorem]{Lemma}
\newtheorem{proposition}[theorem]{Proposition}
\theoremstyle{definition}
\newtheorem{definition}[theorem]{Definition}
\newtheorem{construction}[theorem]{Construction}
\theoremstyle{remark}
\newtheorem{remark}[theorem]{Remark}
\newtheorem{warning}[theorem]{Warning}

\newcommand{\N}{\ensuremath{\mathbb{N}}}
\newcommand{\OT}{\ensuremath{\mathsf{OTree}}}
\newcommand{\WT}{\ensuremath{\mathsf{WTree}}}
\newcommand{\Stage}{\ensuremath{\mathsf{stage}}}
\newcommand{\Reflect}{\ensuremath{\mathsf{reflect}}}
\newcommand{\Index}{\ensuremath{\mathsf{index}}}
\newcommand{\Cseven}{\ensuremath{\mathsf{C}_7}}
\newcommand{\Csix}{\ensuremath{\mathsf{C}_6}}
\newcommand{\lean}[1]{\texttt{#1}}

\title{\textbf{BusyLean Contender 7}\\[0.35em]
\Large Autonomous Critical-Point Reflection over a Self-Generated W-Tree Index}
\author{Research proposal prepared for the BusyLean project}
\date{October 2026}

\begin{document}
\maketitle

\begin{abstract}
The BusyLean repository asks for closed, constructive Lean definitions of very
large natural numbers together with kernel-checked proofs that each new entry
exceeds the preceding one.  The current entry, Contender~6, applies forest
reflection over a finite-type W-language along an index that joins every
$W\N$-tree definable in that same language.  This paper proposes Contender~7,
called \emph{autonomous critical-point reflection}.

The construction first embeds the Contender~6 index into a transparent
Brouwer-tree syntax.  It then closes that seed under ordinal-style addition,
multiplication, exponentiation, fixed-point enumeration, a finite hierarchy of
critical functions, a diagonal common-critical-point operation, and forty-two
levels of higher derivatives.  The resulting autonomous payloads are inserted
into a guarded branch tower.  Every successor branch contains the entire
preceding branch as child~0 and the next autonomous payload as child~1.  At
input~42, the outer reflection root exposes branch~42.  Lean therefore proves a
chain of forty-three strict reflection steps above the complete Contender~6
index, as well as a separate strict lower bound by the level-41 autonomous
payload.

The exact submitted definition is a closed term
\lean{BusyLean.contender\_7 : Nat}.  All recursion is structural; no
\lean{noncomputable}, \lean{partial}, \lean{unsafe}, custom axiom,
\lean{Classical.choice}, previous contender constant, or evaluation shortcut is
used.  The formal theorem is
\[
  \lean{BusyLean.contender\_6\_lt\_contender\_7} :
  \lean{contender\_6 < contender\_7}.
\]
The theorem is independent of any informal ordinal calibration.  The latter is
presented separately and explicitly labelled as heuristic.
\end{abstract}

\tableofcontents

%=====================================================================
\section{Problem and result}

\subsection{The BusyLean setting}

BusyLean is a Lean~4 continuation of the ``name the biggest number'' style of
formal competition.  A legitimate entry is not merely a suggestive notation:
it is a closed Lean definition of type \lean{Nat}, constructive and computable
in principle, accompanied by a kernel-checked strict comparison with the
incumbent.  The repository's engineering rules also exclude the usual escape
hatches: no classical choice in computational dependencies, no custom axioms,
no \lean{sorry}, no \lean{native\_decide} certificate, and no definition that
simply invokes the previous named contender.

The incumbent in the repository at the time of this work is Contender~6:
\begin{lstlisting}
def contender_6 : Nat := stage p5Family index 42
\end{lstlisting}
Here \lean{stage} is forest reflection over a W-type object language, and
\lean{index} is assembled from every $W\N$-tree denoted by an enumerated program
of that language.  Its comparison with the exact Lean port of Contender~5 is
already proved in the repository.

\subsection{The proposed entry}

The new definition is:
\begin{lstlisting}
def contender7Base : Autonomous.OTree :=
  Autonomous.encode index

def contender7RawBranch (k : Nat) : WTree Nat :=
  Autonomous.realize
    (Autonomous.hyper k
      (Autonomous.OTree.succ contender7Base))

def contender7Branch : Nat -> WTree Nat
  | 0 => index
  | k + 1 =>
      .node (fun i =>
        match i with
        | 0 => contender7Branch k
        | _ + 1 => contender7RawBranch k)

def contender7Index : WTree Nat :=
  .node contender7Branch

def contender7Family (n : Nat) : Nat :=
  stage p5Family contender7Index n

def contender_7 : Nat := contender7Family 42
\end{lstlisting}

The construction has two deliberately separate margins.

\begin{enumerate}[leftmargin=*]
\item \textbf{A fully formal guard margin.}  Branch~0 is exactly the complete
  Contender~6 index.  Branch $k+1$ contains branch $k$ as child~0.  Thus, at
  input~42, Lean proves
  \[
    \Stage(B_0,42)<\Stage(B_1,42)<\cdots<\Stage(B_{42},42)
       <\Stage(\mathsf{root},42).
  \]
  This is a chain of forty-three strict reflection steps above the old index.
\item \textbf{An autonomous ordinal-style payload.}  Branch $k+1$ also contains
  as child~1 a realization of the $k$-th higher derivative of a
  common-critical-point closure above the encoded old index.  Hence branch~42
  formally exceeds the level-41 payload, and the final root exceeds branch~42.
\end{enumerate}

The required theorem is proved through the long guard chain rather than by the
shorter direct observation that branch~0 is a child of the root.  This records a
meaningful structural margin in the proof term itself.

%=====================================================================
\section{The inherited reflection engine}

This proposal deliberately reuses the trusted reflection mechanism already
present in BusyLean.  Only the index schedule is strengthened.

\subsection{Oracle families and one-step reflection}

A family is a binary natural-valued oracle:
\[
  \mathsf{Family}=\N\to\N\to\N.
\]
For a family $F$, the existing repository defines
\[
\Reflect(F,n)
 = n+1+
 \max\!\left(
   \max_{i\le n}F(i,n),
   \max_{t\in\mathsf{allTerms}(n)}\mathsf{evalWith}(F,t)
 \right).
\]
The implementation uses primitive list maxima and the enumerator of the
finite-type W-language.  Therefore every visible oracle value and every visible
program value is strictly smaller than \Reflect$(F,n)$.

\subsection{Forest reflection}

For $a:\WT(\N)$, the inherited function \Stage is defined structurally:
\begin{lstlisting}
def stage (seed : Family) : WTree Nat -> Nat -> Nat
  | .leaf   => reflect seed
  | .node f => reflect (fun i => stage seed (f i))
\end{lstlisting}
At a node, the oracle family is precisely the family of recursively computed
child stages.  The key theorem is:
\begin{theorem}[Visible-child domination]
For every seed family $S$, child function $f:\N\to\WT(\N)$, input $n$, and
index $i\le n$,
\[
  \Stage(S,f(i),n)<\Stage(S,\mathsf{node}(f),n).
\]
\end{theorem}

This theorem is the only fact required for the formal comparison.  No
monotonicity of reflection in the oracle is used, because such monotonicity is
false in general: an object-language program may inspect an oracle value and
branch in a way that reverses a pointwise comparison.

\subsection{Contender 6 as the seed schedule}

Contender~6's \lean{index} is a node whose $k$-th child joins, in enumeration
order, all trees denoted by programs of depth at most $k$.  Thus the old
schedule already reflects along the object language's own definable trees.  In
this proposal that entire schedule is treated as one indivisible base object.
It is not re-enumerated, approximated, or replaced by a weaker external
notation.

%=====================================================================
\section{A transparent autonomous tree language}

\subsection{Raw syntax}

We introduce a small Brouwer-tree syntax:
\begin{lstlisting}
inductive OTree : Type where
  | zero  : OTree
  | succ  : OTree -> OTree
  | limit : (Nat -> OTree) -> OTree
\end{lstlisting}
The type is inductive and strictly positive.  Every inhabitant is a finite Lean
value that may contain functions producing further finite values.  Structural
recursion on this W-type is accepted by Lean's termination checker.

One can assign an external ordinal reading by
\[
\begin{aligned}
  |\mathsf{zero}| &= 0,\\
  |\mathsf{succ}(a)| &= |a|+1,\\
  |\mathsf{limit}(f)| &= \sup_{n<\omega}(|f(n)|+1).
\end{aligned}
\]
This interpretation is explanatory only; the submitted Lean comparison does
not depend on an ordinal quotient, ordinal equality, or a formalization of
countable ordinals.

\subsection{Arithmetic constructors}

The source defines syntax-level addition, multiplication, and exponentiation by
structural recursion in the right argument:
\begin{lstlisting}
def add (a : OTree) : OTree -> OTree
  | .zero    => a
  | .succ b  => .succ (add a b)
  | .limit f => .limit (fun n => add a (f n))

def mul (a : OTree) : OTree -> OTree
  | .zero    => .zero
  | .succ b  => add (mul a b) a
  | .limit f => .limit (fun n => mul a (f n))

def pow (a : OTree) : OTree -> OTree
  | .zero    => .succ .zero
  | .succ b  => mul (pow a b) a
  | .limit f => .limit (fun n => pow a (f n))
\end{lstlisting}
With the usual cofinality hypotheses on limit sequences, these are the familiar
recursive clauses for ordinal addition, multiplication, and exponentiation.
The code intentionally avoids asserting quotient-level algebraic laws that have
not been formalized.

Let \lean{omega} be the limit of finite successor chains, and define
\lean{omegaPow a := pow omega a}.  This supplies the level-zero normal-function
syntax $\alpha\mapsto\omega^\alpha$.

\subsection{Fixed-point enumeration}

For a syntax transformer $F:\OT\to\OT$, define
\begin{lstlisting}
def derivative (f : OTree -> OTree) (a : OTree) : OTree :=
  .limit (fun n => iterate f (n + 1) (.succ a))
\end{lstlisting}

The mathematical motivation is Veblen's fixed-point construction.

\begin{proposition}[Classical fixed-point calculation]
Let $F$ be a normal ordinal function, and let $a$ be an ordinal.  Put
$x_0=a+1$, $x_{n+1}=F(x_n)$, and $\lambda=\sup_{n<\omega}x_n$.  Then
$F(\lambda)=\lambda$, and $\lambda$ is the least fixed point of $F$ not below
$x_0$.
\end{proposition}

\begin{proof}
Continuity gives
$F(\lambda)=F(\sup_n x_n)=\sup_n F(x_n)=\sup_n x_{n+1}=\lambda$.
If $y\ge x_0$ and $F(y)=y$, monotonicity gives $x_n\le y$ for every $n$ by
induction, hence $\lambda\le y$.
\end{proof}

The Lean term \lean{derivative} is total without assuming normality.  Normality
is needed only for the intended ordinal interpretation.

\subsection{Critical hierarchy and autonomous closure}

Define finite critical levels by
\[
\begin{aligned}
 C_0(a) &= \omega^a,\\
 C_{k+1}(a) &= C_k'(a),
\end{aligned}
\]
where the prime denotes the preceding fixed-point enumerator.  In Lean:
\begin{lstlisting}
def critical : Nat -> OTree -> OTree
  | 0, a     => omegaPow a
  | k + 1, a => derivative (fun b => critical k b) a
\end{lstlisting}

Next take a diagonal common-critical-point closure:
\begin{lstlisting}
def autonomous (a : OTree) : OTree :=
  .limit (fun k => critical (k + 1) (.succ a))
\end{lstlisting}
The word ``autonomous'' refers to the classical proof-theoretic idea of closing
an ordinal notation under successively stronger critical operations rather than
choosing one fixed external level.  The construction is Veblen-inspired and
Gamma-like, but this paper intentionally does \emph{not} identify its exact
quotient-level order type with $\Gamma_0$ or any other named ordinal.  Such an
identity would require a formal semantic development and a careful analysis of
fundamental sequences.

Finally, iterate the fixed-point operation on the autonomous transformer:
\begin{lstlisting}
def hyper : Nat -> OTree -> OTree
  | 0, a     => autonomous a
  | k + 1, a => derivative (fun b => hyper k b) a
\end{lstlisting}
At input 42, the proposed contender exposes \lean{hyper 41}.  This is a finite
closed definition, not an appeal to an unbounded search or a noncomputable
``least undefinable'' object.

%=====================================================================
\section{Embedding the old index and realizing new schedules}

\subsection{Encoding and realization}

The old reflection engine expects \lean{WTree Nat}, whereas the autonomous
operators act on \lean{OTree}.  Two transparent structural maps connect them:
\begin{lstlisting}
def encode : WTree Nat -> OTree
  | .leaf   => .zero
  | .node f => .limit (fun n => encode (f n))

def realize : OTree -> WTree Nat
  | .zero   => .leaf
  | .succ a =>
      .node (fun n =>
        match n with
        | 0     => realize a
        | _ + 1 => .leaf)
  | .limit f => .node (fun n => realize (f n))
\end{lstlisting}

Under the rank convention
$\operatorname{rk}(\mathsf{leaf})=0$ and
$\operatorname{rk}(\mathsf{node}(f))=\sup_n(\operatorname{rk}(f(n))+1)$,
\lean{realize} has the expected behavior: zero becomes rank~0, successor adds a
visible predecessor child, and limit preserves a countable supremum.  Again,
rank preservation is an intended semantic invariant, not an assumption used by
the strict numerical comparison.

\subsection{Why the seed is not merely an external notation}

The base term is
\[
  B=\lean{Autonomous.encode BusyLean.index}.
\]
Thus the autonomous hierarchy starts \emph{above the actual self-generated
Contender~6 schedule}.  This matters.  A fixed Veblen notation built from zero
could remain far below the proof-theoretic strength already encoded by the
W-language.  Starting from the old index instead applies the critical-point
operations to that stronger base rather than attempting to catch it from below.

The level-$k$ payload is
\[
  R_k=\mathsf{realize}\bigl(\mathsf{hyper}_k(\mathsf{succ}(B))\bigr).
\]
The successor ensures that the first iteration begins strictly above the base
under the intended ordinal reading.

%=====================================================================
\section{The guarded branch tower}

\subsection{Definition}

Let $B_0=\Index$.  Recursively define
\[
  B_{k+1}=\mathsf{node}(g_k),\qquad
  g_k(0)=B_k,\qquad g_k(i+1)=R_k.
\]
This is exactly \lean{contender7Branch}.  The final index is
\[
  J_7=\mathsf{node}(k\mapsto B_k).
\]
At input $n$, the reflection root is guaranteed to see every branch $B_i$ with
$i\le n$.  Therefore input~42 sees $B_{42}$.

\subsection{Why the guard is useful}

The autonomous payload is intended to be much stronger than the old index, but
a proof of that statement at the level of ordinal semantics would be a separate
formalization project.  The guard tower isolates the concerns:

\begin{itemize}[leftmargin=*]
\item the exact strict comparison with Contender~6 follows from the already
  trusted visible-child theorem;
\item the autonomous payload can be studied, strengthened, or replaced without
  endangering that comparison;
\item the proof term explicitly records forty-two strict internal stages rather
  than hiding the old index in a single opaque maximum;
\item child~1 witnesses that the large payload is genuinely present in the
  computed reflection family, not just discussed in comments.
\end{itemize}

The repetition of $R_k$ at every positive child index is intentional.  At any
input $n\ge1$, the payload is visible.  It also avoids an auxiliary default
branch and keeps the definitional reductions at indices 0 and 1 transparent.

%=====================================================================
\section{Formal comparison in Lean}

\subsection{One guarded step}

The repository theorem \lean{stage\_child\_lt} immediately gives:
\begin{lstlisting}
theorem contender7_branch_step_lt (k n : Nat) :
    stage p5Family (contender7Branch k) n <
      stage p5Family (contender7Branch (k + 1)) n := by
  -- child 0 of branch (k+1) is branch k
  exact stage_child_lt p5Family childFamily
    (i := 0) (n := n) (Nat.zero_le n)
\end{lstlisting}
The submitted source spells out \lean{childFamily} as the defining match, so
there is no extensionality argument or hidden rewrite.

Similarly, for $n\ge1$, child~1 yields
\[
  \Stage(R_k,n)<\Stage(B_{k+1},n).
\]

\subsection{The forty-two-step chain}

Primitive recursion on $k$ composes the strict inequalities:
\begin{lstlisting}
theorem contender7_branch_zero_lt_succ :
    forall k n : Nat,
      stage p5Family (contender7Branch 0) n <
        stage p5Family (contender7Branch (k + 1)) n
  | 0, n => contender7_branch_step_lt 0 n
  | k + 1, n =>
      Nat.lt_trans (contender7_branch_zero_lt_succ k n)
        (contender7_branch_step_lt (k + 1) n)
\end{lstlisting}
Instantiating $k=41$ and $n=42$ gives
\[
  \Stage(B_0,42)<\Stage(B_{42},42).
\]
Because $B_0$ is definitionally the old \lean{index}, the left side is exactly
Contender~6.

\subsection{The outer step and final theorem}

Branch~42 is child~42 of the outer root and $42\le42$, so
\[
  \Stage(B_{42},42)<\Stage(J_7,42)=\Cseven.
\]
Transitivity yields:
\begin{theorem}[Kernel-checked comparison]
\[
  \lean{BusyLean.contender\_6 < BusyLean.contender\_7}.
\]
\end{theorem}
The Lean declaration is
\lean{BusyLean.contender\_6\_lt\_contender\_7}.

A second theorem composes this with the repository's existing comparison from
Contender~5:
\[
  \lean{BusyLean.contender\_5 < BusyLean.contender\_7}.
\]

\subsection{Payload inclusion theorem}

Since $R_{41}$ is child~1 of $B_{42}$ and $1\le42$, Lean also proves
\[
  \Stage(R_{41},42)<\Stage(B_{42},42)<\Cseven.
\]
The final declaration is \lean{BusyLean.hyper41\_lt\_contender\_7}.  This
separates the formal fact that the payload participates in the number from the
informal claim about its ordinal magnitude.

%=====================================================================
\section{Computability and constructive status}

\subsection{Why the number is computable in principle}

Every definition used by Contender~7 is a total Lean function:

\begin{itemize}[leftmargin=*]
\item \lean{iterate}, \lean{ofNat}, \lean{critical}, \lean{hyper}, and
  \lean{contender7Branch} recurse structurally on \lean{Nat};
\item \lean{add}, \lean{mul}, \lean{pow}, and \lean{realize} recurse
  structurally on \lean{OTree};
\item \lean{encode} and the inherited \lean{stage} recurse structurally on
  well-founded W-trees;
\item all maxima range over finite lists generated from a finite depth;
\item the final term is closed and has type \lean{Nat}.
\end{itemize}

Lean's acceptance of the definitions is itself a termination certificate for
these recursive equations.  In particular, the construction uses no halting
oracle, Busy Beaver function, minimization over an undecidable predicate, or
choice of an unspecified witness.

\subsection{Why non-evaluation is not noncomputability}

The numerical value is far beyond feasible evaluation.  That is expected and
does not conflict with computability.  A closed structurally recursive program
can have an astronomically long reduction sequence.  The contest compares such
terms intensionally, through proved inequalities, rather than attempting to
print their decimal expansions.

\subsection{Dependency discipline}

The audit checks that \lean{contender\_7} does not depend on the named constants
\lean{contender\_5}, \lean{contender\_6}, their comparison theorems, the old
\lean{maxBy} winner, or P5's named contender constants.  It does reuse the
structural ingredients \lean{index}, \lean{stage}, and \lean{p5Family}; this is
analogous to extending a hierarchy from its previous notation system rather
than inserting the previous final numeral as a seed.

As with Contender~6, the inherited P5 seed family contains its historical
\lean{oldMaximum} guard as a function of depth.  The package discloses this
rather than presenting the dependency closure as cleaner than it is.

%=====================================================================
\section{Growth analysis: formal facts versus heuristic calibration}

\subsection{What is formally established}

The following statements require no ordinal analysis:

\begin{enumerate}[leftmargin=*]
\item Contender~6 is strictly smaller than Contender~7.
\item The proof factors through forty-two guarded branch steps and one outer
  root step.
\item The level-41 autonomous payload is strictly smaller than Contender~7.
\item The definition is total and constructive in Lean.
\end{enumerate}

These are the claims on which adoption as the next BusyLean contender can rest.

\subsection{Intended ordinal reading}

For an informal scale comparison, write $\alpha$ for the rank of the old
self-generated index and $L$ for the effective strength of one reflection step
using the W-language.  The repository's current heuristic describes
Contender~6 as roughly a fast-growing value at an index of the form
$L\cdot\alpha$.

Let $H_0$ denote the diagonal common-critical-point transformer generated by
\lean{autonomous}, and let $H_{k+1}$ enumerate fixed points of $H_k$.  Under the
normal-function interpretation of the OTree constructors, the raw payload
$R_{41}$ has rank on the scale
\[
  H_{41}(\alpha+1).
\]
Consequently the new schedule has rank at least that large, plus the finite
guard and root structure, and the reflected numerical value is heuristically on
a scale resembling
\[
  F_{L\cdot H_{41}(\alpha+1)}(42),
\]
whereas the old entry is on the scale $F_{L\cdot\alpha}(42)$.

This is a qualitative statement: repeated fixed-point enumeration above
$\alpha$ should dwarf multiplication by $\alpha$ or any finite stack of ordinary
reflection steps.

\begin{warning}
The expression above is not a theorem of the submitted Lean development.  It
depends on an external ordinal semantics for OTree, rank preservation of
\lean{encode}/\lean{realize}, normality of the generated transformers, and a
calibration of the object-language reflection operator.  The package therefore
uses words such as ``intended'', ``heuristic'', and ``on the scale of'' rather
than attaching a false precision to a named proof-theoretic ordinal.
\end{warning}

\subsection{Why this still constitutes a significant advance}

Even without the heuristic, the advance is not merely the definition
\lean{contender\_6 + 1}.  The new number performs fresh forest reflection over:

\begin{itemize}[leftmargin=*]
\item the complete old self-generated index;
\item forty-two recursively guarded extensions of that index;
\item an explicit hierarchy of critical-point payloads built above the old
  index, culminating in the level-41 higher derivative;
\item an outer root that exposes the whole branch family to the relativized
  W-language.
\end{itemize}

Thus the construction adds a new schedule language and a new layer of
reflection, while preserving a short, auditable comparison proof.

%=====================================================================
\section{Verification protocol}

The archive includes the exact Lean sources, the workflow used for verification,
and captured logs.  The intended checks are:

\begin{enumerate}[leftmargin=*]
\item cleanly build \lean{BusyLean.Contender7}, which contains definitions only;
\item build the full library including \lean{BusyLean.Comparison7};
\item run \lean{test/Contender7Audit.lean};
\item print the axiom dependencies of the final definition and principal
  theorems;
\item assert that the definition's transitive dependency closure avoids the
  forbidden previous contender constants and winner selectors;
\item record the exact Lean and Lake versions.
\end{enumerate}

The final verification run succeeded on GitHub Actions at repository commit
\begin{center}
\lean{c01efa921d20fcb1d731329ef5582fd16c933fcb}
\end{center}
using Lean~4.34.0 (kernel commit
\lean{293d5d0c0c3f3dded4688b3ccd6a33939ac5102b}) and Lake~5.0.0 on
Ubuntu~24.04.  With both Lake build directories removed, the definition phase
took 1.45 seconds against a 15-second budget, and the remaining proof phase took
2.15 seconds against a 60-second budget.  Seventeen declarations across the
historical and new audits were checked; every reported axiom belonged to the
allowlist $\{\lean{propext},\lean{Quot.sound}\}$.  The new definition and each
principal Contender~7 theorem reported \lean{propext} only.  The transitive
closure of \lean{contender\_7} contained 575 constants and avoided all eight
forbidden previous-winner names.  The feasible sanity check evaluated
\lean{contender7Family 0} to 4; the final input 42 was not evaluated.

The exact checker output and toolchain record are included in the archive and
should be treated as authoritative over any prose in this paper.

%=====================================================================
\section{Limitations and open obligations}

\begin{enumerate}[leftmargin=*]
\item \textbf{No formal ordinal quotient.}  OTree is raw syntax.  Equivalent
  fundamental sequences are not identified, and no well-order comparison
  decision procedure is claimed.
\item \textbf{No formal normality proof.}  The fixed-point interpretation of
  \lean{derivative}, \lean{critical}, \lean{autonomous}, and \lean{hyper} is
  mathematically standard in shape, but normality and exact semantic equations
  are not mechanized here.
\item \textbf{No exact proof-theoretic calibration.}  The fast-growing hierarchy
  labels are explanatory estimates, not kernel-checked bounds.
\item \textbf{Inherited seed complexity.}  The proposal inherits P5's seed
  family and its historical guard.  A future entry could replace this with a
  cleaner independently calibrated family.
\item \textbf{Finite visible level.}  Input~42 exposes \lean{hyper 41}.  This is
  deliberate for a closed contest entry, but an autonomous diagonal over the
  entire higher-derivative sequence may offer a stronger next step.
\end{enumerate}

These limitations do not affect the theorem $\Csix<\Cseven$ or the
computability of \Cseven.

%=====================================================================
\section{Further research questions}

\begin{enumerate}[leftmargin=*]
\item \textbf{Formal rank semantics.}  Define a suitable ordinal or well-founded
  rank semantics for \lean{OTree} and \lean{WTree Nat}; prove that
  \lean{realize} preserves rank and characterize \lean{encode}.
\item \textbf{Normality and derivatives.}  State a syntactic invariant ensuring
  that an OTree transformer represents a normal function, then prove that
  \lean{omegaPow}, \lean{critical}, \lean{autonomous}, and \lean{hyper}
  preserve it.
\item \textbf{Certified lower bounds.}  Formalize a fast-growing hierarchy over
  OTree and prove that the numerical \lean{stage} operation dominates selected
  hierarchy levels.  This would turn the present ordinal heuristic into an
  explicit Lean lower-bound theorem.
\item \textbf{Autonomous diagonalization.}  Replace the fixed index 41 by a
  self-generated schedule that diagonalizes over all higher derivatives while
  remaining structurally recursive and auditable.
\item \textbf{Stronger object languages.}  Implement normalization and bounded
  maxima for System~F, System~$F_\omega$, or a carefully controlled dependent
  calculus, then relativize those languages to the present family.  A stronger
  base language should dominate any improvement that changes only the external
  tree schedule.
\item \textbf{Independent kernel checking.}  Run Lean's independent checker or
  another trusted checker on the compiled environment, and incorporate an
  automated axiom allowlist audit in the repository CI.
\item \textbf{Definition-budget optimization.}  Measure clean builds on the
  official contest hardware, reduce imports and generated code, and keep the
  definition/proof split within the repository's time budgets.
\item \textbf{Eliminate the historical guard.}  Construct a seed family whose
  domination of earlier entries is proved through an independent inclusion
  theorem rather than an inherited maximum over the old language.
\item \textbf{General guard combinator.}  Abstract the guarded-branch pattern and
  prove generic theorems for finite chains, payload visibility, and quantitative
  additive gaps at each reflection step.
\item \textbf{Semantic robustness under alternative fundamental sequences.}
  Determine which growth comparisons survive changes in the realization of
  successor and limit trees, and isolate the invariants actually needed by the
  reflection engine.
\end{enumerate}

%=====================================================================
\section{Integration notes}

The archive is arranged so that its \lean{lean/} directory can be overlaid on a
checkout of BusyLean.  The principal files are:

\begin{center}
\begin{tabular}{ll}
\toprule
File & Purpose\\
\midrule
\lean{BusyLean/Contender7.lean} & definitions only\\
\lean{BusyLean/Comparison7.lean} & strict comparisons and payload theorems\\
\lean{test/Contender7Audit.lean} & axiom and dependency audit\\
\lean{BusyLean.lean} & exports the new modules\\
\lean{.github/workflows/contender7-ci.yml} & reproducible CI verification\\
\bottomrule
\end{tabular}
\end{center}

The development branch used for this proposal is
\begin{center}
\lean{research/autonomous-reflection-c7-20261012}.
\end{center}
Adoption should proceed by reviewing the source and verification logs, not by
trusting the informal size calibration.

%=====================================================================
\section{Conclusion}

Contender~7 combines a conservative proof strategy with an aggressive growth
strategy.  Conservatively, the old Contender~6 schedule is exposed as branch~0,
and forty-two explicit guard nodes plus the outer root provide a transparent
strict comparison.  Aggressively, the same construction carries a hierarchy of
Veblen-inspired critical-point payloads generated above the old self-referential
index, culminating at the forty-first higher derivative visible at input~42.

The resulting Lean term is closed, structurally recursive, and computable in
principle.  Its essential comparison theorem is small enough to audit by hand,
while the autonomous payload supplies a credible and extensible route far
beyond the incumbent's schedule.  The proposal therefore meets the practical
standard for a new BusyLean contender while clearly separating formal theorem,
engineering verification, and proof-theoretic heuristic.

%=====================================================================
\appendix
\section{Principal theorem inventory}

\begin{longtable}{p{0.42\textwidth}p{0.51\textwidth}}
\toprule
Declaration & Meaning\\
\midrule
\endfirsthead
\toprule
Declaration & Meaning\\
\midrule
\endhead
\lean{contender7\_branch\_step\_lt} & Every guarded branch is strictly below its successor at every input.\\
\lean{contender7\_raw\_lt\_next} & Payload $R_k$ is strictly below guarded branch $B_{k+1}$ whenever input is at least one.\\
\lean{contender7\_branch\_zero\_lt\_succ} & The finite transitive closure from branch zero to branch $k+1$.\\
\lean{finalFamily\_lt\_contender7Branch42} & Contender~6's family at input 42 is below branch 42 after forty-two guard steps.\\
\lean{contender7Branch42\_lt\_contender\_7} & The outer root strictly exceeds branch 42.\\
\lean{contender\_6\_lt\_contender\_7} & Required incumbent comparison.\\
\lean{contender\_5\_lt\_contender\_7} & Transitive comparison with the exact Contender~5 port.\\
\lean{hyper41\_lt\_contender7Branch42} & The level-41 autonomous payload is present below branch 42.\\
\lean{hyper41\_lt\_contender\_7} & The final number strictly exceeds the level-41 payload.\\
\bottomrule
\end{longtable}

\section{Reproduction commands}

From a BusyLean checkout containing the patch:
\begin{verbatim}
lake build BusyLean.Contender7
lake build
lake env lean test/Contender7Audit.lean
\end{verbatim}
Do not evaluate \lean{contender\_7}.  The audit evaluates only a feasible
input-zero sanity check for the family.

\section{Bibliographic note}

The use of continuous increasing ordinal functions and fixed-point enumeration
is inspired by Veblen's original hierarchy.  The term ``autonomous'' is used in
the broad proof-theoretic tradition associated with transfinite progressions and
predicative analysis.  The present construction is not claimed to reproduce any
one historical notation system exactly.

\begin{thebibliography}{9}

\bibitem{Veblen1908}
O. Veblen,
\emph{Continuous Increasing Functions of Finite and Transfinite Ordinals},
Transactions of the American Mathematical Society 9 (1908), 280--292.
DOI: \url{https://doi.org/10.1090/S0002-9947-1908-1500814-9}.

\bibitem{Feferman1962}
S. Feferman,
\emph{Transfinite Recursive Progressions of Axiomatic Theories},
Journal of Symbolic Logic 27 (1962), 259--316.

\bibitem{Feferman1964}
S. Feferman,
\emph{Systems of Predicative Analysis},
Journal of Symbolic Logic 29 (1964), 1--30.
DOI: \url{https://doi.org/10.2307/2269764}.

\bibitem{Schuette1965}
K. Schuette,
\emph{Predicative Well-Orderings},
in Formal Systems and Recursive Functions, Studies in Logic and the
Foundations of Mathematics 40 (1965), 280--303.
DOI: \url{https://doi.org/10.1016/S0049-237X(08)71694-X}.

\bibitem{LeanRecursion}
The Lean Project,
\emph{Recursion - Recursive Functions in Lean 4},
\url{https://lean4.dev/language/control-flow/recursion}.

\bibitem{BusyLean}
V. Reshetnikov et al.,
\emph{BusyLean repository},
\url{https://github.com/VladimirReshetnikov/BusyLean}.

\end{thebibliography}

\end{document}
